% ===================================================================
%  કોષ્ટક નં. ૧૧ — શહેર અને તેનાં સ્મારક. --- આગ.
%
%  Traduction, faite en 2026, en gujarati d'aujourd'hui, sur l'ido de
%  texto/io. Regles de la colonne : voir 10-kostak-01.tex. Deux
%  parties, deux numerotations. Ouverture de la TROISIEME SERIE :
%  « ત્રીજી શ્રેણી ».
%
%  DEUX BLOCS MARQUES « ¶2 » PAR ordo.py, ET C'EST LE SEUL TABLEAU DU
%  LIVRET QUI EN AIT DEUX : t11-c1-12-1 et t11-c2-03-1. Une seule cle,
%  DEUX alineas dans chacun. On n'ouvre pas de seconde cle : une ligne
%  blanche dans le bloc, et c'est tout. Sur cinq blocs marques dans
%  tout le releve, deux sont ici.
%
%  LE PREFET (15, 91) EST « કલેક્ટર », ET C'EST LE SEUL RENVOI
%  ADMINISTRATIF DU LIVRET QUI SE TRADUISE SANS PERTE. Le prefet du
%  departement francais et le Collector du district indien sont le
%  meme homme : nomme d'en haut, pose sur une circonscription
%  dessinee a la regle, charge de l'ordre, de l'impot et des
%  catastrophes. Les deux charges ont ete inventees dans le meme
%  siecle par le meme genre d'Etat, l'une a Paris et l'autre a
%  Calcutta, et le mot indien est reste apres le depart de ceux qui
%  l'avaient apporte. « જિલ્લો » rend le departement pour la meme
%  raison. Le tableau 10 disait qu'une langue emprunte les mots de ce
%  qu'on lui a fait ; en voici un qu'elle a garde.
%
%  ET « બંબો » REVIENT, A L'ENDROIT OU ON L'ATTENDAIT LE MOINS. Le mot
%  portugais bomba avait ete releve plus haut comme la pompe de la
%  cour ; ici il est la POMPE A INCENDIE (67, 68), et le pompier est
%  « બંબાવાળો », l'homme du bombo. Avec « ઇસ્પિતાલ » l'hopital (39),
%  qui est hospital, cela fait SEPT mots portugais dans onze tableaux
%  — મેજ la table, પાદરી le pretre, અનનાસ l'ananas, મિસ્ત્રી le
%  macon, બંબો la pompe, ઇસ્પિતાલ l'hopital. Les Portugais ont tenu
%  Diu et Daman quatre siecles et demi, jusqu'en 1961 : ce n'est pas
%  un vieux fonds, c'est un voisinage.
%
%  LE DERNIER MOT DE L'ORCHESTRE EST LE PLUS VIEUX. Cinq instruments
%  a l'alinea 6 : violon (81), flute (82), violoncelle (83),
%  trombone (84), tambour (85). Trois sont transcrits. Deux ont un nom
%  d'ici et ce sont les deux qu'on peut fabriquer soi-meme :
%  « વાંસળી », la flute de bambou — celle de Krishna, qu'on ne peut
%  pas nommer en gujarati sans le nommer lui — et « ઢોલ », le tambour.
%  C'est la loi du tableau 9 verifiee sur la musique : on garde le nom
%  de ce qu'on a manie.
%
%  ET LES COMEDIENS (88, 89) SONT « નટ » ET « નટી », deux mots
%  sanskrits qui nommaient l'acteur et l'actrice quand le theatre
%  indien avait deja ses traites de mise en scene. Ce sont eux qui
%  fuient l'incendie — d'un opera, qu'il a fallu transcrire. LA SCENE
%  EST PLUS VIEILLE ICI QU'EN EUROPE, ET LE SPECTACLE PLUS JEUNE.
%
%  Enfin la porte (8) : « દરવાજો » n'est pas un mot d'archeologie. A
%  Ahmedabad, Delhi Darwaja et Teen Darwaja sont des adresses, on y
%  descend de l'autobus. Le releve grave une porte de remparts
%  medievaux dont il ne reste que l'arc ; la colonne ecrit le mot par
%  lequel des gens se donnent rendez-vous aujourd'hui.
% ===================================================================

%%K t11-apar-1 apar
\VUsaut{2.0mm}
\VUcentre{13.2pt}{90}{ત્રીજી શ્રેણી}
\VUsaut{3.0mm}
\VUfilet{16mm}
\VUsaut{5.6mm}
\VUtitre{15.0pt}{60}{કોષ્ટક નં. ૧૧}
\VUsaut{1.4mm}
\VUpk{11.4pt}{40}{શહેર અને તેનાં સ્મારક. --- આગ.}
\VUsaut{2.2mm}

%%K t11-tit-1 sub
\VUpk{10.2pt}{40}{પરું.}
\VUsaut{3.0mm}

%%K t11-c1-01-1 p
1. --- હું એક મોટા શહેરમાં રહું છું, જે મહાસાગરથી ૧૦૦ કિલોમીટર દૂર છે
અને છતાં પહેલી હરોળનું \VUgras{બંદર} \textsuperscript{(1)} છે. જે નદી
તેને કાંઠે ઘેરે છે તે ૪૦૦ મીટર પહોળી છે અને ખૂબ સુંદર ખાડી બનાવે છે.

%%K t11-c1-02-1 p
2. --- શહેરનો જે આખો ભાગ \VUgras{ઘાટ} \textsuperscript{(2)} પર છે તે
તદ્દન દરિયાઈ અને ઔદ્યોગિક લાગે છે, અને એક \VUgras{પરું}
\textsuperscript{(3)} બનાવે છે, જ્યાં માત્ર ખલાસીઓ અને મજૂરો રહે છે.
તેઓ \VUgras{કારખાનાંમાં} \textsuperscript{(4)} અને \VUgras{ગૅસના
કારખાનામાં} \textsuperscript{(5)} કામ કરે છે, જેનું ઊંચું
\VUgras{ભૂંગળું} \textsuperscript{(6)} અને \VUgras{ગૅસ-ટાંકી} ખૂબ
દૂરથી દેખાય છે.

%%K t11-c1-03-1 p
3. --- મધ્યયુગમાં શહેર કિલ્લેબંધ હતું, અને તેના કોટના થોડા અવશેષ હજી
પણ મળે છે — ખાસ કરીને \VUgras{દરવાજો} \textsuperscript{(8)}, જે જૂના
\VUgras{કિલ્લેબંધ મહેલનો} \textsuperscript{(7)} છે, અને જેની બે બાજુ બે
\VUgras{બુરજ} \textsuperscript{(9)} છે, જે અત્યારે \VUgras{કેદખાના}
તરીકે વપરાય છે.

%%K t11-c1-04-1 p
4. --- આ આખા \VUgras{લત્તામાં}, જે ખૂબ જૂનો લાગે છે, માત્ર એક જ મકાન
પ્રમાણમાં આધુનિક છે : તે \VUgras{જકાતઘર} \textsuperscript{(10)}. તે
ઘાટ અને \VUgras{ગલી} \textsuperscript{(11)} વચ્ચે આવેલું છે, જે જૂના
\VUgras{નગરભવન} \textsuperscript{(12)} સુધી લઈ જાય છે. એ છેલ્લા
સ્મારકને સહેલાઈથી ઓળખી શકાય છે, તેના મહાકાય \VUgras{ઘંટમિનારાને}
\textsuperscript{(13)} લીધે, જે પ્રાચીન શહેર પર પ્રભુત્વ ધરાવે છે.

%%K t11-c1-05-1 p
5. --- રસ્તાની બીજી બાજુએ જકાતઘર જેવી જ શૈલીમાં બાંધેલું એક મકાન
\VUgras{ઊભું છે} : તે \VUgras{શેરબજાર} \textsuperscript{(14)}. તેનો એક
અગ્રભાગ એક નાના ચોક તરફ જુએ છે, જ્યાં \VUgras{કલેક્ટર કચેરી}
\textsuperscript{(15)} છે. ખરું જોતાં અમારું શહેર રાજધાની ન હોવા છતાં
જિલ્લાનું મહત્ત્વનું મુખ્ય મથક છે.

%%K t11-tit-2 sub
\VUsaut{5.0mm}
\VUpk{10.2pt}{40}{શહેરનો સૌથી ઊંચો ભાગ.}
\VUsaut{3.4mm}

%%K t11-c1-06-1 p
6. --- છાવણી, જે સારી એવી મોટી છે, વિશાળ અને ખૂબ આરોગ્યપ્રદ
\VUgras{છાવણી-મકાનોમાં} \textsuperscript{(16)} રહે છે ; એ મકાનો
\VUgras{નગરના બાગની} \textsuperscript{(17)} જાળી સુધી ફેલાયેલાં છે. જે
\VUgras{પહોળો રસ્તો} \textsuperscript{(19)} એ બાગ સુધી લઈ જાય છે તે
\VUgras{બચત બૅન્કની} \textsuperscript{(18)} પાછળથી પસાર થાય છે અને
\VUgras{મુખ્ય દેવળના} \textsuperscript{(20)} ચોકમાં જઈ મળે છે.

%%K t11-c1-07-1 p
7. --- એ બધાં સ્મારક એક નાની ટેકરી પર અર્ધગોળાકારે ફેલાયેલાં છે, જ્યાં
અમારું વિશાળ \VUgras{વિદ્યાલય} \textsuperscript{(21)} પણ છે.

%%K t11-c1-07-2 p
જો કંઈક ઢાળવાળા રસ્તે નીચે ઊતરીએ, જે મોટા રસ્તાને જઈ મળે છે, તો
\VUgras{ન્યાયાલય} \textsuperscript{(22)} આગળ પહોંચાય છે, જેની સામે એક
સુખદ \VUgras{જાહેર બગીચો} \textsuperscript{(26)} ફેલાયેલો છે. એ
\VUgras{બગીચાનો ચોક} બહુ મોટો નથી, પણ તે શહેરની વચ્ચે લીલોતરી અને
ઠંડકનો રણદ્વીપ બનાવે છે. એટલે જ ત્યાં શહેરીઓ ખૂબ આવે છે, જેમને
\VUgras{કૉફીઘરના} \textsuperscript{(25)} છાપરા નીચે બેસવાનું ગમે છે,
લશ્કરના વાજિંત્રીઓને સાંભળવા, જે ગુરુવારે અને રવિવારે
\VUgras{સંગીતમંડપમાં} \textsuperscript{(27)} વગાડે છે, જે લૉન અને
ફૂલક્યારા વચ્ચે મુકાયેલો છે.

%%K t11-c1-08-1 p
8. --- એ લૉનમાંના એક પર કાંસાનું \VUgras{પૂતળું} \textsuperscript{(28)}
ઊભું છે, અમારા એક શહેરીની યાદમાં ઊભું કરાયેલું સ્મારક, જેણે પોતાના
જન્મના શહેરની મોટી સેવા કરી હતી.

%%K t11-tit-3 sub
\VUsaut{4.0mm}
\VUpk{10.2pt}{40}{વાહનવ્યવહાર.}
\VUsaut{3.4mm}

%%K t11-c1-09-1 p
9. --- અત્યાર સુધી અમારું શહેર વીજળીના અજવાળે અજવાળાયું નથી, તેની પાસે
માત્ર \VUgras{ગૅસના દીવા} \textsuperscript{(29)} છે. પણ \VUgras{અહીંનાં
છાપાં} ઝુંબેશ ચલાવે છે કે અજવાળાની એ પદ્ધતિ સુધારાય, \VUgras{જેમ}
\VUgras{ટ્રામ ખેંચવાની પદ્ધતિ} પહેલેથી સુધારાઈ ચૂકી છે.
\VUgras{ઘોડાથી} \VUgras{ખેંચાતી} \VUgras{જૂની} ગાડીઓની જગ્યાએ
વ્યવહારમાં \VUgras{વીજળીની ટ્રામ} \textsuperscript{(31)} આવી ગઈ છે, જે
મુસાફરોને \VUgras{શહેરના એક છેડેથી બીજા છેડે} ખૂબ \VUgras{સસ્તા} ભાવે
ઝડપથી લઈ જાય છે.

%%K t11-c1-10-1 p
10. --- ન્યાયાલય પાસે \VUgras{બૅન્ક} \textsuperscript{(32)} છે, જ્યાં
વેપારી હૂંડીઓ વટાવાય છે. \VUgras{બૅન્કના ઉઘરાણીદાર}
\textsuperscript{(33)} ઘેર જઈને લેણું વસૂલ કરે છે.

%%K t11-tit-4 sub
\VUsaut{4.0mm}
\VUpk{10.2pt}{40}{કળા અને કેળવણી.}
\VUsaut{3.4mm}

%%K t11-c1-11-1 p
11. --- પાસે જે સુંદર મકાન ઊભું છે તે \VUgras{સંગ્રહાલય} છે. તેમાં
એકસાથે શહેરનું \VUgras{પુસ્તકાલય} \textsuperscript{(34)},
\VUgras{કુદરતી વિજ્ઞાનનો સુંદર સંગ્રહ} \textsuperscript{(35)}
(પ્રકૃતિસંગ્રહાલય) અને એક નમણું \VUgras{ચિત્રસંગ્રહાલય}
\textsuperscript{(36)} છે, જે એક ભવ્ય કાયમી \VUgras{પ્રદર્શન} બનાવે છે.

%%K t11-c1-11-2 p
તે અઠવાડિયે બે વાર લોકો માટે ખૂલે છે, પણ બહારગામના લોકો
\VUgras{ચોકીદારને} \textsuperscript{(37)} થોડી બક્ષિસ આપીને ગમે તે
દિવસે તેને જોઈ શકે છે. \VUgras{ચિત્રકારો} \textsuperscript{(38)}
ત્યાં કામ કરવા આવે છે. તેમને પોતાની ઘોડીની \VUgras{આગળ} જોઈ શકાય છે,
હાથમાં \VUgras{રંગપાટી} લઈને, પ્રખ્યાત ચિત્રોની નકલ કરતા.

%%K t11-c1-12-1 p
12. --- ઊંચાણ પર, સંગ્રહાલયની પાછળ, \VUgras{ઘુમ્મટ}
\textsuperscript{(40)} દેખાવા લાગે છે, જે \VUgras{ઇસ્પિતાલનો}
\textsuperscript{(39)} છે, અને \VUgras{અધ્યાપન મંદિરનાં}
\textsuperscript{(41)} મકાનો પણ, જ્યાં અમારી નગરશાળાઓના શિક્ષકોને
તાલીમ અપાતી હતી.

નાટકઘરના ચોકમાં \VUgras{ટપાલ કચેરી} \textsuperscript{(43)} છે, જે એક
\VUgras{ખાનગી મકાનમાં} \textsuperscript{(42)} બેસાડેલી છે.

%%K t11-tit-5 sub
\VUsaut{5.0mm}
\VUpk{11.4pt}{40}{આગ.}
\VUsaut{3.0mm}

%%K t11-tit-6 sub
\VUpk{10.2pt}{40}{આગની બૂમ.}
\VUsaut{3.4mm}

%%K t11-c2-01-1 p
1. --- શહેરમાં એક ભયાનક વાત ફેલાય છે : નાટકઘરમાં હમણાં જ આગ લાગી ! હું
\VUgras{ચોકમાં} \textsuperscript{(96)} પહોંચું છું એ ઘડીએ જ ટોળું
\VUgras{ફૂટપાથ} \textsuperscript{(46)} પર અને \VUgras{પાકા રસ્તા}
\textsuperscript{(47)} પર ભેગું થઈ ગયું છે. \VUgras{ટપાલખાતાના
કારકુનો} \textsuperscript{(44)} અને \VUgras{ટપાલીઓ}
\textsuperscript{(45)} ટપાલ કચેરીમાંથી બહાર નીકળે છે. એક
\VUgras{ટ્રામચાલકે} \textsuperscript{(48)} પોતાની ટ્રામ થોભાવવી પડી,
નગરની \VUgras{એમ્બ્યુલન્સને} \textsuperscript{(50)} જવા દેવા, જે એક
\VUgras{સર્જન} \textsuperscript{(52)} અને એક \VUgras{નર્સ}
\textsuperscript{(51)} સાથે આવે છે, એક \VUgras{ઘાયલની}
\textsuperscript{(53)} સારવાર કરવા, જેને \VUgras{સ્ટ્રેચર}
\textsuperscript{(54)} પર \VUgras{છાપાંના ગલ્લા} \textsuperscript{(55)}
પાસે લવાયો છે.

%%K t11-c2-02-1 p
2. --- પાયદળની એક ટુકડી, જે કસરતથી પાછી ફરતી હતી, \VUgras{ઘોડેસવાર
શહેર-પોલીસ} \textsuperscript{(61)} સાથે વ્યવસ્થા જાળવવા થોભી ગઈ.
\VUgras{ઘોડેસવારો}, જેમના ઘોડા બે પગે ઊભા થઈ જાય છે,
\VUgras{સિપાહીઓ} \textsuperscript{(57)} કે \VUgras{પોલીસ}
\textsuperscript{(63)} કરતાં વધારે સારી રીતે \VUgras{તમાશો જોનારાઓને}
\textsuperscript{(56)} પાછા ધકેલે છે. બાકી પોલીસ ખૂબ કામમાં છે.
તેમાંનો એક \VUgras{હવાલદારોને} \textsuperscript{(64)} બતાવે છે કે બીજા
\VUgras{અકસ્માતના ભોગ બનેલાને} \textsuperscript{(65)} ક્યાં લઈ જવો —
એક બહાદુર બંબાવાળો, જે નિસરણી પરથી પડ્યો છે.

%%K t11-c2-03-1 p
3. --- \VUgras{અફસર} \textsuperscript{(66)} ચોક્કસ જગ્યા બતાવે છે
જ્યાં આવી રહેલો \VUgras{વરાળનો બંબો} \textsuperscript{(67)} ગોઠવવો.
એ શક્તિશાળી યંત્રની રાહ જોતાં જોતાં તેણે તાબડતોબ \VUgras{હાથબંબો}
\textsuperscript{(68)} ગોઠવાવ્યો, જેની લાંબી \VUgras{નળી}
\textsuperscript{(69)} આગ પર ધોધની જેમ પાણી ફેંકે છે.

છ બંબાવાળા તેને ચલાવે છે, જ્યારે તેમનો સાથી, જે \VUgras{ધારાનળી}
\textsuperscript{(70)} પકડે છે, \VUgras{ધારને} \textsuperscript{(71)}
આગની વચ્ચોવચ તાકે છે.

%%K t11-c2-04-1 p
4. --- નાટકઘરની બીજી બાજુએ મોટી \VUgras{નિસરણી} \textsuperscript{(72)}
ટેકવવામાં આવી છે. તેનાથી બંબાવાળા એ જ્વાળાઓની પાસે પહોંચી શકે છે, જે
કાળા \VUgras{ધુમાડાના} \textsuperscript{(75)} વમળ વચ્ચેથી ફૂટે છે.

%%K t11-tit-7 sub
\VUsaut{5.0mm}
\VUpk{10.2pt}{40}{બહાદુરીનાં કામ.}
\VUsaut{3.4mm}

%%K t11-c2-05-1 p
5. --- \VUgras{આગ} \textsuperscript{(74)} \VUgras{નાટકઘરના}
\textsuperscript{(73)} પ્રવેશખંડમાં જન્મી, જ્યારે \VUgras{ઑપેરાની}
તાલીમ ચાલતી હતી, જેનો પહેલો \VUgras{ખેલ} બીજા દિવસે થવાનો હતો.
સદ્ભાગ્યે પ્રેક્ષાગૃહમાં થોડા જ \VUgras{પ્રેક્ષકો}
\textsuperscript{(76)} હતા. બિચારા \VUgras{ઝરૂખા}
\textsuperscript{(77)} પર આશરો લેવા ગયા, જ્યાં બહાદુર
\VUgras{બંબાવાળા} \textsuperscript{(86)} નિસરણી અને દોરડાં લઈને તેમને
બચાવવા આવે છે.

%%K t11-c2-06-1 p
6. --- \VUgras{થાંભલાવાળી પરસાળ} \textsuperscript{(78)} નીચે, જ્યાં
\VUgras{જાહેરાતનું પાટિયું} \textsuperscript{(79)} ખેલની ખબર આપે છે,
ત્યાં \VUgras{વાદકો} \textsuperscript{(81)} દેખાય છે, જે
\VUgras{વાદ્યવૃંદના} \textsuperscript{(80)} છે અને જે પોતાનાં વાજિંત્રો
લઈને ભાગે છે : \VUgras{વાયોલિન} \textsuperscript{(81)},
\VUgras{વાંસળી} \textsuperscript{(82)}, \VUgras{ચેલો}
\textsuperscript{(83)}, \VUgras{ટ્રોમ્બોન} \textsuperscript{(84)},
\VUgras{ઢોલ} \textsuperscript{(85)}, વગેરે. \VUgras{રાચરચીલાનો}
\textsuperscript{(87)} થોડો ભાગ બચાવવાની કોશિશ થાય છે.

%%K t11-c2-07-1 p
7. --- \VUgras{નટ} \textsuperscript{(89)} અને \VUgras{નટીઓ}
\textsuperscript{(88)}, \VUgras{ગાયકો} અને \VUgras{ગાયિકાઓને} પોતાના
નાટકના \VUgras{પોશાક} \textsuperscript{(90)} સાથે જ ભાગવું પડ્યું.
ટેનર એક \VUgras{પત્રકારને} \textsuperscript{(92)} સમજાવે છે કે એ કેવી
રીતે બન્યું. \VUgras{ખબરપત્રી} પહેલી જ ઘડીથી અધિકારીઓ સાથે દોડી આવ્યો
હતો : \VUgras{કલેક્ટર} \textsuperscript{(91)}, \VUgras{મેયર}
\textsuperscript{(93)}, \VUgras{પોલીસ કમિશનર} \textsuperscript{(94)}
અને \VUgras{ન્યાયાધિકારી} \textsuperscript{(95)}, જે જવાબદારી નક્કી
કરવા તપાસ કરશે.

\VUsaut{6.0mm}
\VUfilet{22mm}
